\chapter{Signal Chain and Theory of Operation}
\label{ch:signal}

This chapter follows the signal from digital synthesis, out through the DAC and
analog front end, back in through the ADC, through quadrature demodulation to
magnitude/phase, and into the packetizer and trigger engine. Module names refer
to \file{FPGA/Custom Hardware Source Files/}; the demodulator cores live inside
the packaged \code{modulator\_chain} IP.

\section{Clock and reset domains}
Two synchronous domains, both from the \SI{200}{\mega\hertz} oscillator:
\begin{itemize}[nosep]
  \item \textbf{\SI{200}{\mega\hertz} (\code{clk\_out1})} --- synthesis and
    demodulation front end (DDS, mixers, \code{dac\_multitone\_mixer},
    \code{am\_modulator}, \code{demod\_frame\_combiner}).
  \item \textbf{\SI{100}{\mega\hertz} (\code{clk\_out2})} --- control/output:
    \code{output\_packetizer\_v2}, \code{timestamp\_counter}, the AXI register
    file, DMA, and the \code{trigger\_engine}.
\end{itemize}
The \SI{200}{}$\to$\SI{100}{\mega\hertz} crossing is a single AXIS CDC FIFO
between the frame combiner and the packetizer. \code{reset\_manager} generates
the per-domain synchronized resets (using \code{reset\_sync}) from
\reg{SYS\_CTRL0}. The trigger fire path stays entirely in the
\SI{100}{\mega\hertz} timestamp domain, so the output edge has no CDC.

\section{Excitation synthesis (DDS $\to$ DAC)}
Each of four \code{modulator\_chain} instances contains a rasterized DDS
(Xilinx DDS Compiler, full-range fix16.15) producing a cosine carrier and its
quadrature sine. \code{dac\_multitone\_mixer} scales each cosine by its per-tone
amplitude (\reg{CH$n$\_FA[31:16]}, unsigned Q1.15, \code{0x8000} = unity), gates
it by the per-tone enable (\reg{CHANNEL\_EN[3:0]}), sums the four, and saturates
to 16 bits. The sine outputs are \emph{not} summed --- they are the internal
quadrature local oscillator for demodulation. The frequency of tone $n$ is
$f_n = \text{FCW}_n \times \SI{12.5}{\kilo\hertz}$ (Section~\ref{sec:fpga-regs}).
The composite goes through \code{dac\_output\_wrapper} (which serializes the
16-bit samples to 8 DDR LVDS pairs with a FRAME strobe) to the AD9122.

To keep the multi-tone sum from clipping, the host splits unity across the
enabled tones (1 tone = full scale, 2 = half each, 4 = quarter each); the mixer's
saturation is a hard backstop.

\section{Analog path}
DAC $\to$ reconstruction/buffer amplifier $\to$ SMA $\to$ sensor electrodes; the
returning current $\to$ low-noise transimpedance front end $\to$ programmable-gain
amplifier (\reg{ANALOG1} sets its gain) $\to$ differential driver $\to$ AD9467.
Components are detailed in Chapter~\ref{ch:hardware}. The excitation must sit
above the \SI{400}{\kilo\hertz} analog high-pass cutoff.

\section{Acquisition and input routing}
\code{adc\_lvds\_capture} deserializes the AD9467 LVDS output into 16-bit samples
in the \SI{200}{\mega\hertz} domain. A \code{mux} (\code{mux.v}, selected by
\reg{SYS\_CTRL0[6]}) chooses the demodulator input: the live ADC, or an internal
loopback from \code{am\_modulator}. The loopback path AM-modulates a replayed
message onto the DDS carrier (envelope $1 + m\,x(t)$, $m = 0.5$), providing a
self-test signal with no analog hardware. For replay, \code{sample\_hold}
paces PC-pushed samples out at ${\sim}625$\,kSa/s
($\SI{200}{\mega\hertz}/320$).

\section{Quadrature demodulation (the lock-in)}
\label{sec:lockin}
Each \code{modulator\_chain} multiplies the input by the cosine and sine of its
reference to form the in-phase and quadrature products, low-pass filters and
decimates them (CIC + FIR), and converts the complex result to polar form with a
CORDIC:
\[
X = x[n]\cos(\omega_\text{ref} n), \qquad
Y = x[n]\sin(\omega_\text{ref} n),
\]
\[
R = \sqrt{X^2 + Y^2}, \qquad \varphi = \operatorname{atan2}(Y, X).
\]
The CORDIC polar output arrives on \code{r\_phase\_out\_tdata[63:0]} as
\textbf{magnitude \code{fix32\_30}} in \code{[31:0]} and \textbf{phase
\code{fix32\_29}} in \code{[63:32]}. In engineering units:
\[
\text{mag} = \frac{\text{sig}}{2^{30}}, \qquad \varphi = \frac{\text{phase}}{2^{29}}\ \text{rad}.
\]
The demod output rate is $\sim$\SI{200}{\kilo\hertz} per channel.

\section{Framing and packetization}
\code{demod\_frame\_combiner} takes the four chains' polar buses directly and
packs them, the raw-ADC monitor, and the per-channel over-range (OR) flags into
one 288-bit (36-byte) ``super-record'' per demod instant:
\[
\{\underbrace{\text{OR}[1:0], 14'b0, \text{adc}[15:0]}_{\text{status}},\ r_3,\ r_2,\ r_1,\ r_0\},
\]
and hands it to the \SI{200}{}$\to$\SI{100}{\mega\hertz} CDC FIFO. In the
\SI{100}{\mega\hertz} domain, \code{output\_packetizer\_v2} stamps each record
with the local timestamp and serializes \emph{only} the channels selected by
\code{channel\_mask} into 32-bit AXIS beats (field order in
Section~\ref{sec:datastream}). \code{records\_per\_packet} sets \code{TLAST} =
one S2MM DMA buffer; both mask and count are latched at packet boundaries so a
mid-stream change cannot corrupt framing. This is the self-describing format the
PC deinterleaves. The stream then goes over AXI DMA (S2MM) to the PS.

\section{Timebase and the trigger engine}
\code{timestamp\_counter} is a free-running 64-bit counter clocked at
\SI{100}{\mega\hertz} but incrementing by 2, so one count is \SI{5}{\nano\second}.
Its value stamps every record, so every sample the PC sees carries an absolute
board time.

\code{trigger\_engine} (Stage 2, dual output) turns a PS decision into a
jitter-free edge. The PS loads a target timestamp $T$ and pulse width $W$ per
output into spare AXI registers, then pulses a shared fire strobe; each output
goes HIGH exactly when the counter reaches its $T$, holds for $W$, and returns
LOW. Because $T$ is read straight from the same counter that stamped the data,
there is no clock synchronization anywhere. The two outputs (main AA7,
secondary AA6) implement the reference's dual-actuator pattern via
$T_1 = T_0 + \text{pair\_delay}$; a width of 0 disables an output.

\paragraph{Detection (PS).} The stateful detection --- rolling-baseline bipolar
peak detection, per-channel phase-window classification, and fixed/velocity delay
math --- runs in \code{iza\_replay.c} on the PS, not the fabric
(\file{guides/firmware-guide.md}, \file{guides/triggering-guide.md}).

\section{Module catalog}
The active custom modules and their role in the chain:
\begin{center}
\begin{longtable}{@{}l l@{}}
\toprule
Module & Role \\
\midrule\endhead
\code{modulator\_chain} (IP) & per-tone DDS + quadrature demod (CIC/FIR/CORDIC) \\
\code{dac\_multitone\_mixer} & sum/scale/gate the four carriers into the DAC excitation \\
\code{dac\_output\_wrapper}  & serialize 16-bit samples to AD9122 LVDS (byte mode) \\
\code{am\_modulator}         & AM loopback / self-test signal \\
\code{sample\_hold}          & pace replayed samples at \SI{625}{\kilo\hertz} \\
\code{mux}                   & demod input select (ADC vs loopback) \\
\code{adc\_lvds\_capture}, \code{adc\_bit\_reorder}, \code{adc\_fifo\_read} & AD9467 LVDS capture path \\
\code{demod\_frame\_combiner}& pack 4 channels + ADC + OR into a super-record \\
\code{output\_packetizer\_v2}& self-describing packetizer to DMA \\
\code{timestamp\_counter}    & 5 ns free-running timebase \\
\code{trigger\_engine}       & dual fire-at-timestamp outputs \\
\code{reset\_manager}, \code{reset\_sync} & per-domain synchronized resets \\
\code{spi\_rd\_tristate}     & SPI read tristate for the converters \\
\code{AXI\_register\_file} (IP) & the 16 control registers \\
\bottomrule
\end{longtable}
\end{center}
Retired modules (superseded, moved to \file{retired hardware/}) are listed in the
FPGA developer guide.
